\documentclass[10pt,a4paper]{amsart}
\usepackage[margin=19mm]{geometry}
\usepackage{amsmath,amssymb,mathtools}
\usepackage[scr=rsfs,cal=euler]{mathalfa}
\usepackage{xfrac}
\usepackage[unicode,hidelinks,bookmarks=false]{hyperref}
\newcommand{\ie}{i.e.\ }
\newcommand{\cf}{cf.\ }
\newcommand{\resp}{respectively\ }
\newcommand{\Subsection}{\subsection}
\providecommand{\nobreakdash}{\nobreak\hbox{-}}
\setlength{\emergencystretch}{2em}
\multlinegap0pt
\usepackage{bm}
\usepackage{bbm}
\usepackage{dsfont}
\usepackage{xspace}
\usepackage{cancel}
\usepackage{mathtools}
\usepackage{amsthm}
\makeatletter
\@ifundefined{example}{\newtheorem{example}{Example}}{}
\@ifundefined{lemma}{\newtheorem{lemma}{Lemma}}{}
\@ifundefined{remark}{\newtheorem{remark}{Remark}}{}
\makeatother
\title[The most symmetric smooth cubic surface]{The most symmetric smooth cubic surface}
\author{Anastasia V. Vikulova}
\date{Source: arXiv:2407.09997v3 (CC BY 4.0)}
\begin{document}
\maketitle

\noindent\emph{Source and licence.} The most symmetric smooth cubic surface. Version arXiv:2407.09997v3 (CC BY 4.0), distributed under the Creative Commons Attribution 4.0 International licence, as recorded in the arXiv licence metadata for this exact version and shown on the abstract page https://arxiv.org/abs/2407.09997v3. Full rights notice: licenses/arxiv-2407.09997-CC-BY-4.0.txt.\par\smallskip

\noindent\textit{Editorial scope.} Counted unit: the Lemma whose source label is \texttt{l53} in the article (source0.tex lines 2170--2176, source label \texttt{l53}), delivered together with the article's own complete original proof of it (source0.tex lines 2178--2269, one \texttt{proof} environment of 92 lines). No mathematics is written, repaired or re-derived: statement and proof are the article's own text line for line. Required context retained uncounted: l50 (lines 2026--2029); r6 (lines 2055--2058); e4 (lines 2060--2079); eq7.4 (lines 2063--2077). Every definition, display and label of those lines is retained unchanged. Omitted from this package and not redistributed: 1--2025, 2030--2054, 2059--2059, 2078--2169, 2177--2177, 2270--2555. Nothing inside a retained excerpt is omitted or altered: every retained body is delivered exactly as the article's lines are written, so no omission receipt of any kind accompanies this package. Citations: the retained excerpts carry no citation command at all, so no bibliography entry of the article is transcribed and no reference is redistributed. Reference adaptation: none was needed and none was made; every reference inside the delivered unit resolves inside it, so no unresolved reference or citation is produced. Printed slips kept literal: no printed word, symbol, hypothesis or number of the counted statement or of its proof was altered. Layout and class: the article's own class is \texttt{amsart} with packages including amsmath, amssymb, array, amscd, leftindex, tikz-cd, float, calc, color, ulem; the wrapper is the lane's pinned \texttt{amsart} editorial wrapper at 10pt on A4, which restates the theorem-like environments the retained text needs and adds the article's own macro definitions that the retained text needs (none), with extra wrapper packages (bm, bbm, dsfont, xspace, cancel, mathtools). The delivered numbering is the unit's own. Assets: no retained span contains a figure environment, a graphics-inclusion command, a PDF-page inclusion, a table or a TikZ picture; no image file is copied into this package, the package declares no asset dependency and no image file is redistributed with it. Licence: version arXiv:2407.09997v3 of this article is distributed under the Creative Commons Attribution 4.0 International licence, as recorded in the arXiv licence metadata for this exact version and shown on the abstract page https://arxiv.org/abs/2407.09997v3. A full rights notice is delivered at licenses/arxiv-2407.09997-CC-BY-4.0.txt. Characters: every retained excerpt is pure ASCII, so the delivered engine, XeLaTeX with its default Latin Modern text font, reports no missing character.
\begin{lemma}
\label{l50}
Up to conjugation the group $(\mathbb{Z}/3\mathbb{Z})^3 \rtimes \mathfrak{S}_4$ contains a~unique subgroup isomorphic to $(\mathbb{Z}/3\mathbb{Z})^2 \rtimes \mathfrak{D}_4$X.
\end{lemma}

\begin{remark}
\label{r6}
From now on we mean by $(\mathbb{Z}/3\mathbb{Z})^2 \rtimes \mathfrak{D}_4$ a~subgroup of $(\mathbb{Z}/3\mathbb{Z})^3 \rtimes \mathfrak{S}_4$, which is unique up to conjugation by Lemma~\ref{l50}.
\end{remark}

\begin{example}
\label{e4}
Let $\mathbf{F}$ be a~field not containing nontrivial cube roots of unity. Consider the projective space $\mathbb{P}^3$ with homogeneous coordinates $x$, $y$, $z$ and~$t$. Consider the subgroup isomorphic to $(\mathbb{Z}/3\mathbb{Z})^2$ in the automorphism group of $\mathbb{P}^3$ that is generated by
\begin{equation}
\label{eq7.4}
\begin{pmatrix}
-1 & 1 & 0 & 0\\
-1 & 0 & 0 & 0\\
0 & 0 & 1 & 0\\
0 & 0 & 0 & 1
\end{pmatrix} \quad \text{and} \quad
\begin{pmatrix}
1 & 0 & 0 & 0\\
0 & 1 & 0 & 0\\
0 & 0 & -1 & 1\\
0 & 0 & -1 & 0
\end{pmatrix}.
\end{equation}
Consider the group $\mathfrak{D}_4 \subset \mathrm{PGL}_4(\mathbf{F})$ generated by the permutations~$(12)$ and~$(1324)$ of the homogeneous coordinates $x$, $y$, $z$ and~$t$ in~$\mathbb{P}^3$. Then $(\mathbb{Z}/3\mathbb{Z})^2$ and~$\mathfrak{D}_4$ generate the group $(\mathbb{Z}/3\mathbb{Z})^2 \rtimes \mathfrak{D}_4$ mentioned in Remark~\ref{r6}.
\end{example}

\begin{lemma}
\label{l53}
Let $\mathbf{F}$ be a~field of characteristic different from $2$ and~$3$ such that it contains no nontrivial cube roots of unity. Let $S$ be a~smooth cubic surface over~$\mathbf{F}$ such that its automorphism group is isomorphic to $(\mathbb{Z}/3\mathbb{Z})^2 \rtimes \mathfrak{D}_4$. Then $S$ is isomorphic to the cubic surface
$$
2(x^3+y^3+z^3+t^3)-3(x^2y+xy^2+z^2t+zt^2)=0.
$$
\end{lemma}

\begin{proof}
First of all, recall that by Lemma~\ref{l50} and Remark~\ref{r6} the group
$$
(\mathbb{Z}/3\mathbb{Z})^2 \rtimes \mathfrak{D}_4
$$
is isomorphic to the group defined in Example~\ref{e4}.

Consider all cubic surfaces invariant under the action of $(\mathbb{Z}/3\mathbb{Z})^2$ generated by the elements~\eqref{eq7.4}. These cubic surfaces are defined by homogeneous polynomials~$f(x,y,z,t)$ of degree~$3$ such that each element $g \in (\mathbb{Z}/3\mathbb{Z})^2$ preserves $f(x,y,z,t)$ up to a~scalar factor. Since the matrices~\eqref{eq7.4} are of order~$3$ and $\mathbf{F}$ does not contain nontrivial cube roots of unity, this scalar is equal to~$1$. Thus, first let us find all homogeneous polynomials of degree~$3$ invariant under the action of
\begin{equation}
\label{eq7.7}
\begin{pmatrix}
-1 & 1 & 0 & 0\\
-1 & 0 & 0 & 0\\
0 & 0 & 1 & 0\\
0 & 0 & 0 & 1
\end{pmatrix}.
\end{equation}

Consider a~homogeneous polynomial
\begin{align}
\notag
f(x,y,z,t)
&=ax^3+by^3+cx^2y+dxy^2
\\
&\qquad+x^2l_1(z,t)+y^2l_2(z,t)+xyl_3(z,t)+xg_1(z,t)+yg_2(z,t)+h(z,t),
\label{eq7.8}
\end{align}
where $a,b,c,d \in \mathbf{F}$, $l_1(z,t)$, $l_2(z,t)$ and $l_3(z,t)$ are linear homogeneous polynomials,~$g_1(z,t)$ and~$g_2(z,t)$ are quadratic homogeneous polynomials, and $h(z,t)$ is a~cubic homogeneous polynomial. The element~\eqref{eq7.7} maps $f(x,y,z,t)$ to
\begin{align*}
\widetilde{f}(x,y,z,t)
&=(-a-b-c-d)x^3+ay^3+(3a+2c+d)x^2y-(3a+c)xy^2
\\
&\qquad
+(l_1(z,t)+l_2(z,t)+l_3(z,t))x^2+l_1(z,t)y^2{-}\,(2l_1(z,t)+l_3(z,t))xy
\\
&\qquad
-(g_1(z,t)+g_2(z,t))x+g_1(z,t)y+h(z,t).
\end{align*}
It is not difficult to see that a~homogeneous polynomial $f(x,y,z,t)$ such that
$$
f(x,y,z,t)=\widetilde{f}(x,y,z,t)
$$
is of the form
$$
f(x,y,z,t)=a(x^3+y^3-3xy^2)+bxy(x-y)+l(z,t)(x^2+y^2-xy)+h(z,t),
$$
where $a,b \in \mathbf{F}$, $l(z,t)$ is a~linear homogeneous polynomial, and $h(z,t)$ is a~cubic homogeneous polynomial. Applying the same reasoning to the matrix
$$
\begin{pmatrix}
1 & 0 & 0 & 0\\
0 & 1 & 0 & 0\\
0 & 0 & -1 & 1\\
0 & 0 & -1 & 0
\end{pmatrix}
$$
gives us a~homogeneous polynomial
\begin{equation}
\label{eq7.9}
f(x,y,z,t)=a(x^3+y^3-3xy^2)+bxy(x-y)+c(z^3+t^3-3zt^2)+dzt(z-t),
\end{equation}
where $a,b,c,d \in \mathbf{F}$. So all cubic surfaces invariant under the action of~$(\mathbb{Z}/3\mathbb{Z})^2$ generated by the elements~\eqref{eq7.4} are of the form~\eqref{eq7.9}. Among the surface defined by polynomials~\eqref{eq7.9} let us find the cubic surfaces that are also semi-invariant under the action of the group $\mathfrak{D}_4$ generated by the permutations $(12)$ and $(1324)$ of~the coordinates $x$, $y$, $z$ and~$t$. Consider the permutation $(12)$. It takes~\eqref{eq7.9} to the~homogeneous polynomial
\begin{equation}
\label{eq7.10}
\widehat{f}(x,y,z,t)=a(x^3+y^3-3yx^2)-bxy(x-y)+c(z^3+t^3-3zt^2)+dzt(z-t).
\end{equation}
We can see from~\eqref{eq7.10} that if $f(x,y,z,t)=-\widehat{f}(x,y,z,t)$, then $f(x,y,z,t)=0$ defines a~singular cubic surface. This means that the polynomials~\eqref{eq7.9} invariant under the action of the permutation $(12)$ and defining smooth cubic surfaces are of the form
\begin{equation}
\label{eq7.11}
f(x,y,z,t)=a(2(x^3+y^3)-3(x^2y+xy^2))+c(z^3+t^3-3zt^2)+dzt(z-t),
\end{equation}
where $a \in \mathbf{F}^*$ and $c,d \in \mathbf{F}$ are such that~\eqref{eq7.11} defines a~smooth surface. The same reasoning, as applied to the permutation
$$
(12)(1324)^2=(34),
$$
shows that smooth cubic surfaces must be defined by polynomials of the form
\begin{equation}
\label{eq7.12}
f(x,y,z,t)=a(2(x^3+y^3)-3(x^2y+xy^2))+c(2(z^3+t^3)-3(z^2t+zt^2)),
\end{equation}
where $a,c \in \mathbf{F}^*$. Using the semi-invariance of~\eqref{eq7.12} under the permutation $(1324)$ we obtain
\begin{equation}
\label{eq7.13}
f(x,y,z,t)=2(x^3+y^3)- 3(x^2y+xy^2)\pm (2(z^3+t^3)-3(z^2t+zt^2)).
\end{equation}
However, the change of coordinates
$$
x \mapsto x, \qquad y \mapsto y, \qquad z \mapsto -z \quad\text{and}\quad t \mapsto -t
$$
shows that the smooth cubic surfaces defined by the polynomials~\eqref{eq7.13} with plus and minus signs are isomorphic.

Lemma~\ref{l53} is proved.
\end{proof}

\end{document}
